\section{Proofs for the boundary results}\label{app:lim}

This appendix proves \cref{prop:lim:rotating,thm:lim:posslp,prop:lim:flow,prop:lim:rank}.

\subsection{Proof of \texorpdfstring{\cref{prop:lim:rotating}}{the rotating-fiber proposition}}\label{app:lim:rotating}

Order the coordinates as $(v,z_1,z_2)$ and put $r=z_1-vz_2$. Then
$\nabla r=(-z_2,1,-v)$, and the only nonzero second derivatives of $r$ are
$\partial_{vz_2}r=\partial_{z_2v}r=-1$. Write $E$ for the symmetric matrix
with entries $-1$ in positions $(v,z_2)$ and $(z_2,v)$ and zeros elsewhere.

\emph{Part \ref{it:lim:rot-core}.} For fixed $v$, the residual term $r^2$
vanishes at $z=0$, so $V_\gamma(v)=(v-\frac12)^2+\gamma v$. Completing the
square gives $V_\gamma(v)=(v-a_\gamma)^2+c_\gamma$ with
$a_\gamma=(1-\gamma)/2\in(0,1)$ and a constant $c_\gamma$. The residual
gradient of $r$ is $(1,-v)$ and its residual Hessian is zero, so the
residual Hessian of $F_\gamma$ is $2(1,-v)^{\mathsf T}(1,-v)\succeq0$. Finally
$\partial_{vv}F_\gamma=2+2(\partial_vr)^2=2+2z_2^2\le4$.

\emph{Part \ref{it:lim:rot-fiber}.} By part \ref{it:lim:rot-core},
$F_\gamma-c_\gamma=(v-a_\gamma)^2+r^2$, which vanishes exactly when
$v=a_\gamma$ and $z_1=a_\gamma z_2$. With $z_2=s\in[0,1]$, the point
$z_1=a_\gamma s$ lies in $[0,1]$.

\emph{Part \ref{it:lim:rot-indef}.} The Hessian is
\[
 H=\diag(2,0,0)+2\nabla r\,\nabla r^{\mathsf T}+2rE .
\]
Let $w=(0,v,1)$. Then $\nabla r^{\mathsf T}w=v-v=0$ and $w^{\mathsf T}Ew=0$,
so $w^{\mathsf T}Hw=0$, while $Hw=2rEw=(-2r,0,0)$. A positive semidefinite
matrix with $w^{\mathsf T}Hw=0$ satisfies $Hw=0$, so $H$ is not positive
semidefinite when $r\ne0$. Explicitly, for $d=w+te_v$ with
$t=r/(1+z_2^2)$,
\[
 d^{\mathsf T}Hd=2t\,e_v^{\mathsf T}Hw+t^2H_{vv}
 =-4rt+(2+2z_2^2)t^2=-\frac{2r^2}{1+z_2^2}<0 .
\]
Since $H_{z_1z_1}=2>0$, $H$ is indefinite. The set $\{r=0\}$ is the zero
set of a nonzero polynomial and has empty interior. A twice differentiable
function that is convex on a convex set $K$ with nonempty interior has a
positive semidefinite Hessian at every interior point of $K$, and the
interior of $K$ contains points with $r\ne0$. Hence $F_\gamma$ is convex on
no such $K$. Linear terms do not change $H$, so the conclusion holds for
every $\gamma$.

\emph{Part \ref{it:lim:rot-strict}.} Put $\tilde r=a-vb$. For fixed $v$,
the residual objective $q_v(a,b)=\tilde r^2+b^4$ is strictly convex: on a
segment along which $b$ is not constant, $b^4$ is strictly convex and
$\tilde r^2$ is convex; on a segment with constant $b$, $a$ varies and
$\tilde r^2$ is strictly convex in $a$. Its unique minimizer is $a=b=0$,
independently of $v$, so the projected value, its minimizer $a_\gamma$, and
the growth identity of part \ref{it:lim:rot-core} are unchanged, and the
unique optimizer is $(a_\gamma,\frac12,\frac12)$, which is interior. The
residual Hessian is $2(1,-v)^{\mathsf T}(1,-v)+\diag(0,12b^2)\succeq0$ and
$\partial_{vv}\widetilde F_\gamma=2+2b^2\le4$. At a point with $b=0$ the
quartic term has zero Hessian and $\partial_v\tilde r=-b=0$, so the
computation of part \ref{it:lim:rot-indef} applies with $r$ replaced by
$\tilde r=a$: with $w=(0,v,1)$ in the coordinates $(v,a,b)$ and $d=w+ae_v$,
one gets $d^{\mathsf T}Hd=-4a^2+2a^2=-2a^2<0$ when $a\ne0$. Every
neighborhood of the optimizer contains such points.\qed

\subsection{Proof of \texorpdfstring{\cref{thm:lim:posslp}}{the PosSLP theorem}}\label{app:lim:posslp}

A straight-line program is a sequence of integers $A_1,\ldots,A_s$ in which
each $A_j$ is $0$, $1$, or $A_l\circ A_m$ with $l,m<j$ and
$\circ\in\{+,-,\times\}$; its output is $A=A_s$. The integers can have
exponentially many bits. The construction never computes them; it writes
only bounded local formulas.

\begin{lemma}[Bounded pair encoding]\label{lem:lim:pairs}
Define rational pairs $(u_j,v_j)$ by $(0,\frac14)$ for the constant $0$,
$(\frac14,\frac14)$ for the constant $1$, and, for $A_j=A_l\pm A_m$ and
$A_j=A_lA_m$ respectively,
\[
 u_j=\frac{u_lv_m\pm u_mv_l}4,\quad v_j=\frac{v_lv_m}4,\qquad
 u_j=\frac{u_lu_m}4,\quad v_j=\frac{v_lv_m}4 .
\]
Then $v_j>0$, $u_j=A_jv_j$, and $\abs{u_j},\abs{v_j}\le\frac14$ for every $j$.
The number $t=(2u_s-v_s)/4=v_s(2A-1)/4$ is nonzero, it is positive if and
only if $A>0$, and $\abs t\le\frac3{16}$.
\end{lemma}

\begin{proof}
By induction, $v_j$ is a product of positive numbers, and
$u_j/v_j=u_l/v_l\pm u_m/v_m$ or $(u_l/v_l)(u_m/v_m)$, which is $A_j$. If all
earlier coordinates lie in $[-\frac14,\frac14]$, the new numerators have
absolute value at most $2\cdot\frac1{16}\cdot\frac14=\frac1{32}$ or
$\frac1{64}$. Since $A$ is an integer, $2A-1\ne0$, and $v_s>0$.
\end{proof}

List the coordinates $u_1,v_1,\ldots,u_s,v_s,t$ as $x_1,\ldots,x_N$, $N=2s+1$.
Each is given by a rule $x_i=p_i(x_1,\ldots,x_{i-1})$, where $p_i$ is a
constant, one of the bilinear formulas of \cref{lem:lim:pairs}, or the
affine formula for $t$. On the box $\mathcal B=[-\frac14,\frac14]^N$ these
rules satisfy
\begin{equation}\label{eq:lim:gate-bounds}
 \abs{p_i}\le\tfrac14,\qquad\norm{\nabla p_i}_1\le1,\qquad
 \norm{\nabla^2p_i}_2\le1 .
\end{equation}
Indeed, each bilinear rule has at most four gradient entries of absolute
value at most $\frac1{16}$, or fewer and larger entries when $l=m$, with
$\norm{\nabla p_i}_1\le\frac14$; its Hessian has at most two nonzero entries
in each row, of absolute value at most $\frac14$, or a single entry at most
$\frac12$ when $l=m$; and the rule for $t$ is affine with
$\norm{\nabla p_N}_1=\frac34$. The exact assignment $x^*$, defined by
$x_i^*=p_i(x_{<i}^*)$, lies in $\mathcal B$ by \cref{lem:lim:pairs}, and
$x_N^*=t^*$ has the sign of $A-\frac12$.

\begin{lemma}[Weighted gate objective]\label{lem:lim:gates}
Let $G(x)=\sum_{i=1}^N64^{N-i}\bigl(x_i-p_i(x)\bigr)^2$. Then
$\nabla^2G\succeq I$ on $\mathcal B$, $x^*$ is the unique zero of $G$ in
$\mathcal B$, and $G(x)\ge\frac12\norm{x-x^*}^2$ on $\mathcal B$.
\end{lemma}

\begin{proof}
Write $a_i=64^{N-i}$, $\rho_i=x_i-p_i(x)$, $D=\diag(\sqrt{a_i})$, and let $L$
be the strictly lower-triangular matrix with $L_{ij}=\partial_jp_i$. For a
direction $h$,
\[
 h^{\mathsf T}\nabla^2Gh=2\norm{D(I-L)h}^2-2\sum_ia_i\rho_i\,h^{\mathsf T}\nabla^2p_ih .
\]
Put $\widehat E=DLD^{-1}$, so $\widehat E_{ij}=8^{j-i}L_{ij}$ for $j<i$. By
\eqref{eq:lim:gate-bounds} every row of $\widehat E$ has absolute sum at
most $\frac18$, and every column has absolute sum at most
$\sum_{l\ge1}8^{-l}=\frac17$. Hence
$\norm{\widehat E}_2\le(\frac18\cdot\frac17)^{1/2}<\frac14$ and
\[
 2\norm{D(I-L)h}^2=2\norm{(I-\widehat E)Dh}^2\ge\tfrac98\norm{Dh}^2 .
\]
On $\mathcal B$, $\abs{\rho_i}\le\frac12$, and $p_i$ depends only on
$x_{<i}$, so $\abs{h^{\mathsf T}\nabla^2p_ih}\le\norm{h_{<i}}^2$. Therefore the
second term is at least
\[
 -\sum_ia_i\norm{h_{<i}}^2=-\sum_jh_j^2\sum_{i>j}a_i\ge-\frac1{63}\sum_ja_jh_j^2 .
\]
Altogether $h^{\mathsf T}\nabla^2Gh\ge(\frac98-\frac1{63})\norm{Dh}^2\ge\norm h^2$,
since every $a_j\ge1$. Because the rules are triangular, $G(x)=0$ exactly
when $x=x^*$. All residuals vanish at $x^*$, so $\nabla G(x^*)=0$, and
Taylor's theorem on the segment from $x^*$ to $x$ gives the last claim.
\end{proof}

\begin{lemma}[Two sign tests]\label{lem:lim:signs}
Let $\varepsilon=\frac1{32}$ and, on $\mathcal B\times[0,1]$, let
\[
 H_+(x,a)=G(x)+a^2-\varepsilon ax_N,\qquad H_-(x,a)=G(x)+a^2+\varepsilon ax_N .
\]
Each is strongly convex with modulus $1-\varepsilon$ and has a unique
minimizer $(x^\pm,a_\pm)$. Moreover $a_+>0$ if and only if $t^*>0$, and
$a_->0$ if and only if $t^*<0$. Exactly one of $a_+,a_-$ is positive.
\end{lemma}

\begin{proof}
The Hessian of $G(x)+a^2$ is at least $I$, and the bilinear term has Hessian
of operator norm $\varepsilon$; this gives the modulus and uniqueness on the
compact convex box. On the face $a=0$ both functions equal $G$, whose unique
minimizer on $\mathcal B$ is $x^*$, with $\nabla G(x^*)=0$. At $(x^*,0)$ the
$x$-gradient of $H_\pm$ is $\nabla G(x^*)=0$ and the $a$-derivative is
$\mp\varepsilon t^*$. If this derivative is nonnegative, $(x^*,0)$
satisfies the first-order conditions of the convex problem and is the
minimizer, so $a_\pm=0$. If it is negative, increasing $a$ decreases the
objective, so $(x^*,0)$ is not the minimizer; and no other point of the face
$a=0$ is, because a minimizer on that face would minimize $G$. Hence
$a_\pm>0$. Since $t^*\ne0$, exactly one case gives a positive value.
\end{proof}

\begin{lemma}[A convex amplifier]\label{lem:lim:amplifier}
Let $\eta=\frac1{192}$ and, on $Y=(\mathcal B\times[0,1])^2\times[0,1]$, let
\[
 \Phi(x^+,a_+,x^-,a_-,y)=H_+(x^+,a_+)+H_-(x^-,a_-)
 +\eta\bigl[a_+^2(y-1)^2+a_-^2y^2\bigr].
\]
Then $\Phi$ is convex on $Y$ and has a unique minimizer, whose
$y$-coordinate is $1$ if $t^*>0$ and $0$ if $t^*<0$.
\end{lemma}

\begin{proof}
For constants $c$ and a direction $(p,q)$ in the variables $(a,y)$, direct
differentiation gives
\[
 (p,q)^{\mathsf T}\nabla^2\bigl[a^2(y-c)^2\bigr](p,q)
 =2\bigl(aq+2(y-c)p\bigr)^2-6(y-c)^2p^2 .
\]
With $c=1$ for $a_+$ and $c=0$ for $a_-$, and $\abs{y-c}\le1$, the
amplifier contributes at least $-6\eta(p_+^2+p_-^2)$, where $p_\pm$ are the
components of the direction along $a_\pm$. The Hessian of $H_++H_-$ is at
least $(1-\varepsilon)I$ on all base coordinates. So the quadratic form of
$\nabla^2\Phi$ is at least $(1-\varepsilon-6\eta)\norm{p_{\mathrm{base}}}^2
=\frac{15}{16}\norm{p_{\mathrm{base}}}^2\ge0$, and $\Phi$ is convex.

The amplifier is nonnegative, so $\Phi\ge\min H_++\min H_-$. If $t^*>0$,
\cref{lem:lim:signs} gives $a_+>0=a_-$, and choosing the two base
minimizers and $y=1$ makes the amplifier vanish; so this lower bound is the
minimum. Every minimizer must then attain both base minima, which fixes all
base coordinates, and must make the amplifier vanish, which forces $y=1$
because $a_+>0$. The case $t^*<0$ is symmetric and forces $y=0$.
\end{proof}

\begin{proof}[Proof of \cref{thm:lim:posslp}]
The map writes $\Phi$ of \cref{lem:lim:amplifier}: two copies of the
$N$ local rules, the weights $64^{N-i}$ of $O(N)$ bits, and constant-size
remaining coefficients. Every residual $x_i-p_i(x)$ has at most three
monomials, so $\Phi$ is an explicit rational quartic with $O(N)$ monomials,
computed in time polynomial in the size of the program. By
\cref{lem:lim:amplifier,lem:lim:pairs}, $\Phi$ is convex on the rational box
$Y$, has a unique minimizer $\hat x$, and $\hat x_y=1$ exactly when $A>0$.

For $F_\gamma$, the core and residual terms are separate. The core term
$(v-\frac12)^2+\gamma v$ equals $(v-v_\gamma)^2$ plus a constant, so the
core optimizer is $v_\gamma=(1-\gamma)/2\in[\frac14,\frac34]$, the projected
growth is one, and $\partial_{vv}F_\gamma=2$. The residual objective $\Phi$
does not depend on $\gamma$, so the unique optimizer is
$(v_\gamma,\hat x)$ on every draw.

For (a), given a program, construct $\Phi$, compute the law, draw $\gamma$,
run the algorithm, and compute a point within distance $1/4$ of
$(v_\gamma,\hat x)$. Its $y$-coordinate is within $1/4$ of $\hat x_y$, so
comparing it with $\frac12$ decides whether $A>0$. The answer is correct on
every draw, and the expected running time is polynomial. For (b), apply the
deterministic algorithm to $\Phi$ directly. Square Root Sum is decidable in
polynomial time with a $\PosSLP$ oracle~\citep{AllenderEtAl2009}, so
$\PosSLP\in\mathrm P$ implies that Square Root Sum is in $\mathrm P$.
\end{proof}

\subsection{Proof of \texorpdfstring{\cref{prop:lim:flow}}{the flow-boundary proposition}}\label{app:lim:flow}

The atoms of $U_{1,M}$ are $-1+2j/(M-1)$, $0\le j<M$. Such an atom is at
least $\frac12$ exactly when $j\ge3(M-1)/4$. For a power of two $M\ge4$,
$3M/4$ is an integer and $3M/4-1<3(M-1)/4\le3M/4$, so the condition is
$j\ge3M/4$, which holds for $M/4$ indices. Each coordinate therefore has
probability $\frac14$, and by independence $\Pr(\mathcal E)=\frac1{16}$. The
second derivatives of $F_\gamma$ in the core are $\partial_{v_iv_i}=1$, so
$L=1$ is valid.

\emph{Part \ref{it:lim:flow-growth}.} Minimizing $v_1z_1+v_2z_2$ over the
two labels gives $\min\{v_1,v_2\}$, so $V_\gamma$ has the stated form. On
$\mathcal E$ each $\gamma_i>0$, and $v_i\ge v_i^2$ on $[0,1]$, so
$\gamma^{\mathsf T}v\ge\min_i\gamma_i\norm v^2$; with $\min\{v_1,v_2\}\ge0$
this gives $V_\gamma(v)\ge(\frac12+\min_i\gamma_i)\norm v^2\ge\norm v^2$.
Since $V_\gamma(0)=0$, the origin is the unique optimal core, and both
labels have cost zero there.

\emph{Part \ref{it:lim:flow-box}.} At $(a,0)$ with $a>0$ the label $(1,0)$
costs $a$ and $(0,1)$ costs $0$; at $(0,b)$ with $b>0$ the opposite holds.
The noise term is common to both labels and cancels from these comparisons.
So each label is strictly worse than the other at some point of the box,
and no valid certificate can assert that one label is optimal throughout.

\emph{Part \ref{it:lim:flow-search}.} The cell $[0,h]^2$ has the corner $0$
with value $V_\gamma(0)=\min V_\gamma$, and the incumbent $U_j$ is a
feasible value, hence at least $\min V_\gamma$. So the lower bound of the
cell, at most $V_\gamma(0)-h^2/4$, does not exceed $U_j$, and the cell is
retained whenever its parent is. At level $0$ the whole core box is
retained, and induction gives the claim at every level.

\emph{The chain.} With $m-1$ further stages of two zero-cost parallel arcs,
each feasible flow chooses one arc per stage, so there are $2^m$ flows. The
conditional value is unchanged, so parts
\ref{it:lim:flow-growth}--\ref{it:lim:flow-search} hold, and on
$\mathcal E$ every flow is optimal at the origin. In every box
$[0,a]\times[0,b]$ the optimal flows at $(a,0)$ and at $(0,b)$ use
different first-stage arcs, so no single flow is certified on the retained
hull at any level. A strategy that enumerates all flows when this
certificate fails does at least $2^m$ work on $\mathcal E$, and its
expected work is at least $2^m/16$. The network has $2m$ unit-capacity arcs
and $m+1$ nodes, so its encoding has length $O(m\log m)$.\qed

\subsection{Proof of \texorpdfstring{\cref{prop:lim:rank}}{the rank-separation proposition} and a nonquadratic convexification}\label{app:lim:rank}

\emph{Part \ref{it:lim:rank-convex}.} The residual Hessian is
$\theta t\nabla^2P(z)$ with
$\nabla^2P(z)=12\diag(z_i^2)+12(\sum_iz_i)^2\mathbf1\mathbf1^{\mathsf T}\succeq0$,
which is positive semidefinite for $t\ge0$, and $\partial_{tt}F=\lambda$.
At $t=0$ the Hessian is
\begin{equation}\label{eq:lim:rank-hessian}
 \nabla^2F(0,z)=\begin{pmatrix}\lambda&\theta\nabla P(z)^{\mathsf T}\\
 \theta\nabla P(z)&0\end{pmatrix},
\end{equation}
and for any $i$ with $\partial_iP(z)\ne0$ its principal $2\times2$ minor on
$(t,z_i)$ has determinant $-\theta^2(\partial_iP(z))^2<0$. So $F$ is not
jointly convex.

\emph{Part \ref{it:lim:rank-rank}.} If $F+\frac12(t,z)^{\mathsf T}K(t,z)$ is
convex on the box, then $\nabla^2F+K\succeq0$ at interior points, and by
continuity of the Hessian and closedness of the positive semidefinite cone
also at points with $t=0$. Let $u\in\ker K_{zz}$ and $w=(0,u)$. Then
$w^{\mathsf T}Kw=0$, and since $K\succeq0$, $Kw=0$. By
\eqref{eq:lim:rank-hessian}, $w^{\mathsf T}(\nabla^2F(0,z)+K)w=0$, so the
positive semidefinite matrix $\nabla^2F(0,z)+K$ annihilates $w$. Its first
component gives $\theta\nabla P(z)^{\mathsf T}u=0$ for every
$z\in[0,1]^m$. At $z=e_i$, $\nabla P(e_i)=4(e_i+\mathbf1)$, and the matrix
$4(I+\mathbf1\mathbf1^{\mathsf T})$ with these columns is invertible. Hence
$u=0$, the positive semidefinite block $K_{zz}$ is positive definite, and
$\rank K\ge\rank K_{zz}=m$. Linear terms do not affect any Hessian, so
\ref{it:lim:rank-convex} and \ref{it:lim:rank-rank} hold after adding them.

\emph{Part \ref{it:lim:rank-interior}.} At $t\in\{0,1\}$,
$F\ge\frac\lambda8-\sum_ib_i>0$ because $P\ge0$ and $z\le\mathbf1$. If
$b_i>0$, then $F(\frac12,\epsilon e_i)=\theta\epsilon^4-b_i\epsilon<0$ for
$0<\epsilon<(b_i/\theta)^{1/3}$. So the minimum is negative and is not
attained at $t\in\{0,1\}$. Both inequalities are strict, so they survive
linear perturbations whose $\ell_1$ norm is smaller than half the smaller
of the two margins.\qed

\begin{remark}[A one-coordinate nonquadratic correction]\label{rem:lim:entropy}
Extend $t\log t$ by $0$ at $t=0$. For $C\ge\frac43\theta\max_{[0,1]^m}P$, the
function $F+Ct\log t$ is convex on the box. Indeed, for $t>0$ its Hessian is
\[
 \begin{pmatrix}\lambda+C/t&\theta\nabla P^{\mathsf T}\\
 \theta\nabla P&\theta t\nabla^2P\end{pmatrix}.
\]
Since $P$ is homogeneous of degree four, Euler's identity gives
$\nabla^2P(z)z=3\nabla P(z)$, so $\nabla P$ lies in the range of $\nabla^2P$
and $\nabla P^{\mathsf T}(\nabla^2P)^{+}\nabla P=\frac13\nabla P^{\mathsf T}z
=\frac43P$; the kernel component of $(\nabla^2P)^+\nabla P-z/3$ is
orthogonal to $\nabla P$. The generalized Schur complement is therefore
$\lambda+(C-\frac43\theta P(z))/t\ge0$, and the Hessian is positive
semidefinite for $t>0$. Convexity on $(0,1]\times[0,1]^m$ and continuity
give convexity on the closed box. So
\cref{prop:lim:rank} separates the small core from fixed quadratic
corrections only, not from every difference-of-convex representation. The
scale $C$ grows with $\theta$, and the correction has an unbounded
derivative at $t=0$.
\end{remark}
